\chapter{“Vinte segundos”}
\label{sec:1}
Um oficial de inteligência, ao descrever seu papel numa esteira de seleção de alvos construída em torno de listas de morte geradas por máquina, contou a um repórter que gastava vinte segundos em cada nome, o suficiente para verificar que o alvo era um homem: “Não leva muito tempo para saber se alguém tem voz de homem ou de mulher”. Acrescentou que, como ser humano, seu valor agregado era zero, a não ser como carimbo de aprovação. “Economizava muito tempo” \autocite{abraham2024lavender}. No relato dele, era a isso que se resumia o “humano no circuito” sob o ritmo operacional.\footnote{A pesquisa sobre viés de automação prevê essa delegação: quando outras demandas disputam a atenção delas, as pessoas que recebem a recomendação de uma máquina a verificam menos e a acatam mais \autocite{parasuraman2010complacency}.}

Uma interpretação é que o oficial é um sujeito a quem deram tempo de menos, e que a solução é governá-lo melhor — mais tempo, um segundo revisor, consciência adicional instalada no circuito. Por essa lente da \emph{sujeição social}, a pessoa é um sujeito que está fora de uma máquina e a opera como ferramenta: um motorista e um carro, um operário e um torno. A pessoa tem um papel, um nome, responsabilidades e uma consciência a que se pode dirigir a palavra. Uma lente alternativa é a \emph{servidão maquínica}, em que a pessoa é um componente da máquina, montado com as outras peças num sistema que as regula juntas; a máquina usa a pessoa como um de seus órgãos, e o que a pessoa contribui é uma função executada num ritmo que o sistema define.\footnote{O nome que Doctorow dá à posição do oficial é centauro reverso: não uma pessoa assistida por uma máquina, mas um componente de máquina feito de uma pessoa \autocite{doctorow2026reversecentaur}.} As duas não são mutuamente excludentes; uma pessoa pode ter um papel num plano enquanto é uma peça no outro \autocite[Platô 13, “The Apparatus of Capture”: há servidão quando os próprios seres humanos são peças constituintes de uma máquina que compõem entre si e com outras coisas, sob o controle e a direção de uma unidade superior, e sujeição quando a unidade superior constitui o ser humano como sujeito ligado a um novo objeto exterior, que pode ser um animal, uma ferramenta ou mesmo uma máquina]{deleuze1987plateaus}, e o oficial dos vinte segundos tinha as duas coisas. Como peça, ele era o lugar onde a esteira produzia seu carimbo, e vinte segundos era o tempo de ciclo dessa peça. Como sujeito, ele foi interpelado o tempo todo: descreveu a um jornalista a dissonância entre as notícias de mortes de civis e a pressão no trabalho por mais alvos. E, mesmo quando a dissonância chegou à sua consciência, ele seguiu em frente. Ela se fez notar e não decidiu nada;\footnote{Ou, pelo menos, nada que impedisse as Forças de Defesa de Israel de matar 100 mil pessoas em Gaza entre outubro de 2023 e outubro de 2025 \autocite{gomezugarte2025uncertainty,acosta2026gaza,economist2025gaza}. E o relato desse oficial específico foi publicado em abril de 2024 pela +972 Magazine. Ninguém sabe dizer se a decisão desse oficial de falar com um jornalista evitou uma morte a mais, 50 mil mortes a mais ou nenhuma. E o que teria acontecido se o oficial tivesse se recusado depois de duas semanas na função — digamos, hipoteticamente, em novembro de 2023? Estamos falando de um oficial de inteligência que usava o Lavender para montar listas de morte. O ambiente de trabalho dele no dia a dia provavelmente estava mais para um emprego de escritório civil do que para um filme de ação. Ele deveria ter buscado reuniões repetidas com os superiores, subindo pela cadeia de comando, enquanto se recusava a aprovar alvos para ataques aéreos dos quais se podia razoavelmente esperar mortes de civis? Dá para argumentar que essa teria sido a coisa ética a fazer. Nesse contrafactual, teria havido menos mortos?} o julgamento de que a política era desproporcional veio, disse ele, em retrospecto.

Então o que faltava não era uma consciência; ele tinha uma. Uma consciência funciona fazendo uma decisão custar alguma coisa, mas aqui o saldo dos custos pendia para o outro lado: “Éramos pressionados o tempo todo: ‘Tragam mais alvos’. Eles gritavam mesmo com a gente” \autocite{abraham2024lavender}. Ele era um componente com escrúpulos, mas o ritmo se manteve: os escrúpulos não mudaram nada do que ele contribuía, e ele virou um carimbo. Um componente com escrúpulos que não segura o ritmo vira o “gargalo humano” que um ex-comandante da unidade de inteligência de sinais de Israel escreveu um livro sobre como eliminar \autocite{ys2021humanmachine}. O que fazia do oficial um componente era o arranjo, não a atitude dele. Um componente pode seguir em frente a contragosto ou com entusiasmo: o funcionário irrefletido que Arendt viu no julgamento de Eichmann \autocite{arendt1963eichmann} era, a julgar pelo que ele tinha dito na Argentina quatro anos antes, sem nada em jogo, um número encenado para Jerusalém \autocite{stangneth2014eichmann}. O que a esteira deixava o oficial contribuir era uma operação fixa num ritmo fixo: ele via cada caso como a esteira o descrevia, um de cada vez, no tempo que a esteira lhe dava. Se alguma coisa poderia ocupar essa posição e não ser um componente é a pergunta da parte IV.

Reportagens investigativas descreveram três sistemas funcionando juntos em Gaza: um gerando alvos de edifícios e infraestrutura, com fontes dizendo que a maioria das sugestões eram residências particulares; um segundo pontuando indivíduos pela suposta filiação militante, produzindo listas que teriam marcado até 37 mil pessoas e que, internamente, se sabia estarem erradas cerca de uma vez em cada dez\footnote{Os militares israelenses contestaram esse relato. Os relatos de Minab e do Maven, mais adiante, se baseiam em reportagens investigativas que as partes também contestam. Eu os tomo como foram noticiados.}; e um terceiro rastreando quando as pessoas marcadas voltavam para casa, para que os ataques aéreos fossem programados para quando elas estivessem com a família.

\runin{Maven: a versão americana, agora com aceleração por IA generativa}

O bombardeio da escola em Minab matou 156 pessoas, 120 delas crianças em idade escolar \autocite[par.~92]{UNFFMIran2026}.\footnote{Os números variam conforme a fonte e a data, de 156 (o balanço final do promotor de Minab, também dado pela Anistia Internacional) a 168 (o prefeito da cidade, nos primeiros dias). A missão da ONU dá os números do promotor e também relata informações independentes de 157 mortos, 123 deles crianças de 13 anos ou menos \autocite[pars.~92–93]{UNFFMIran2026}.} A função civil da escola era visível em imagens de satélite de fontes abertas, e a codificação do alvo ainda trazia uma designação herdada de uma instalação vizinha da Guarda Revolucionária. Uma missão de apuração de fatos da ONU relatou haver motivos razoáveis para crer que o bombardeio constituiu o crime de guerra de tornar um bem civil objeto de ataque \autocite{ohchr2026minabpress}, e que a escola era claramente identificável, sem nenhuma informação que indicasse uso militar \autocite[pars.~81–86, 124–126]{UNFFMIran2026}. Concluiu que os Estados Unidos descumpriram a obrigação de fazer tudo o que fosse viável para verificar o alvo, incluindo a atualização dos dados de seleção de alvos entre as precauções viáveis, e situou a falha além da negligência \autocite[pars.~124–125]{UNFFMIran2026}. A Anistia Internacional chegou a uma conclusão paralela, em termos que acompanham o dever de verificar do artigo 57(2)(a)(i) do Protocolo Adicional I: que as forças podem ter se baseado em inteligência desatualizada e deixado de fazer tudo o que fosse viável para verificar o alvo \autocite{Amnesty2026Minab}.

Aconteceu na manhã de 28 de fevereiro, o primeiro dia da campanha do Irã \autocite{Amnesty2026Minab,crs2026iranwar}. Na véspera, o secretário tinha anunciado que o fornecedor (Anthropic) cujo modelo (Claude) ficava dentro da plataforma de seleção de alvos (Maven Smart System) seria designado um risco para a cadeia de suprimentos por se recusar a licenciar o Claude para todo uso legal \autocite{npr2026openai}. O Maven Smart System tinha devolvido a escola como alvo recomendado para o primeiro dia \autocite{livingstone2026maven,gizmodo2026minab}. Nenhum relato publicado diz que papel o Claude teve nessa recomendação, e uma pessoa familiarizada com o trabalho do fornecedor disse à NBC News que o modelo não dava diretamente conselhos sobre alvos \autocite{nbc2026targeting}.

Uma revisão interna do Pentágono, cujas conclusões completas continuam sob sigilo rigoroso, teria constatado que a preparação de alvos que levava horas em conflitos anteriores foi condensada em minutos, que o pessoal de mitigação de danos a civis caiu cerca de noventa por cento, para menos de vinte pessoas, que a equipe do Comando Central passou de dez para um, e que nenhum membro dela revisou o local de Minab — um resultado que as autoridades atribuíram tanto aos cortes quanto à decisão de não envolver a equipe no planejamento e nas operações \autocite{Bloomberg2026KillChain}. A investigação formal tinha sido aberta em março, depois que uma avaliação preliminar concluiu que os Estados Unidos eram provavelmente responsáveis \autocite{NPR2026MinabProbe}.

Na campanha do Irã de 2026, uma plataforma de seleção de alvos construída por uma empresa contratada, com o modelo de propósito geral de um fornecedor comercial dentro dela, usada para classificar alvos por importância estratégica e redigir uma justificativa jurídica para cada ataque \autocite{livingstone2026maven}, ajudou a atingir mais de mil alvos nas primeiras vinte e quatro horas, cerca de dez vezes o que era possível antes dela \autocite{csis2026maven}, e treze mil até o trigésimo oitavo dia \autocite{breakingdefense2026maven}. O Pentágono não publicou nenhum tempo de revisão por alvo para ela.\footnote{Maven é um nome de programa que abriga mais de um sistema: o programa do Pentágono que ainda se chama Maven, e o Maven Smart System, a plataforma de software que nasceu dele \autocite{csis2026maven}. O programa começou em 2017 como um projeto para rotular objetos nos vídeos que drones coletavam durante a campanha contra o Estado Islâmico \autocite{weisgerber2017maven,fang2019maven}. O sistema de 2026 funde sistemas militares e fontes de inteligência numa única interface, usada para encontrar e atingir alvos e para comandar operações \autocite{defensescoop2026transition}. A plataforma é da Palantir: o diretor-chefe digital e de IA do Pentágono a chama de “o produto de software como serviço da Palantir” \autocite{defensescoop2026transition}. A posição pública da Palantir, citada de segunda mão, é que decidir sobre vida ou morte não é o papel dela e que a responsabilidade sempre fica com a organização militar \autocite{lee2026palantir}, e ela sustenta que a inteligência de base era responsabilidade do governo e que não há provas de que o Maven tenha falhado ou causado o bombardeio de Minab \autocite{Bloomberg2026KillChain}. Como o programa passou de um escritório e de uma empresa contratada para os seguintes está no apêndice~\ref{sec:C.1}.}

As pessoas nas listas de ataques aéreos israelenses e americanas não teriam como se opor a estar nessas listas — as esteiras não foram construídas para obter esse tipo de retorno. Andry José Hernández Romero foi apanhado por uma esteira diferente, e ele podia ter se oposto. Maquiador venezuelano, ele se apresentou no posto de fronteira de San Ysidro em agosto de 2024 para pedir asilo, passou na entrevista de medo crível e tinha advogado, um pedido pendente e uma data de audiência \autocite{LULAC2025HernandezRomero,Vega2025HernandezRomero}. Um funcionário de uma empresa contratada que fazia a triagem dele interpretou duas tatuagens, coroas sobre as palavras para “mamãe” e “papai”, como marcas de gangue \autocite{Palomera2025Blade}, e em março de 2025 ele foi levado de avião para uma prisão em El Salvador antes da audiência, na noite em que um juiz federal mandou que os aviões dessem meia-volta \autocite{jgg2025contempt}. Ele tinha tudo de que uma pessoa precisa para se opor, menos o tempo. Voltaremos ao caso dele com mais detalhe quando tratarmos da esteira americana de sequestro e transferência forçada.

\textbf{Quem se recusou.} O registro dos vinte segundos tem uma segunda metade, o registro de quem resistiu a uma pressão desse tipo, e ele diz que a composição da sala importa muito. Na condição de base de Milgram, cerca de um terço dos sujeitos parou antes do choque máximo. Com dois pares desobedecendo, só um em cada dez foi até o fim, e, quando dois experimentadores discordavam, muito menos gente seguiu adiante do que na condição de base \autocite{milgram1974obedience}. Cerca de um quarto dos sujeitos de Asch nunca cedeu à maioria errada \autocite{asch1956studies}. Browning estima que entre um décimo e um quinto dos homens do Batalhão 101 da Polícia de Reserva se recusaram a atirar em Józefów, o primeiro massacre do batalhão, ou se esquivaram \autocite{browning1992ordinary}. A proporção ia de quase ninguém a quase todos, conforme o ambiente. Quem resistiu tinha evidências de fora do canal: um par que desobedecia, uma segunda autoridade que discordava e, entre os salvadores que os Oliner estudaram, relações que atravessavam a linha que a perseguição traçava \autocite{oliner1988altruistic}. Ou mantinham o compromisso com uma precisão que a pressão não conseguia vencer, como se mantêm os mandatos morais e os valores sagrados \autocite{skitka2009limits,ginges2007sacred}. Ou tratavam a fonte como parte interessada. E grupos resistiam onde indivíduos isolados raramente resistiam \autocite{gamson1982encounters}.

A sala também pode ser composta a favor do outro lado. A Alemanha de Weimar tinha uma vida associativa densa, o que possivelmente ajudou a derrubar a república: à medida que partidos e governos deixavam de responder às queixas das pessoas, elas punham sua energia em clubes e sociedades organizados segundo linhas de grupo, e os nazistas recrutaram através das redes que essas sociedades já tinham construído \autocite{berman1997weimar}. Cidades com mais corais, clubes esportivos e associações de ginástica viram uma adesão mais rápida ao partido, mas só onde o governo era instável; na Prússia, cujo governo era estável, a relação era pequena e não significativa \autocite{satyanath2017bowling}. Então o que os salvadores tinham não era simplesmente mais companhia. Os laços deles atravessavam a linha que a perseguição traçava. Ashutosh Varshney encontrou o mesmo em cidades indianas: tumultos entre hindus e muçulmanos eram raros onde as associações reuniam as duas comunidades nas mesmas organizações, e comuns onde não reuniam \autocite{varshney2002ethnic}. Os laços que tornam a recusa possível são os que um movimento fascista teria de cortar.

Tudo o que vem a seguir é sobre compor a sala.\footnote{A sala é o ambiente, e, como diria minha professora de teatro, o cenário dirige a cena. Quem mais está na sala? Dá para sobreviver a sair? Alguém ouviria uma recusa?} Compor é uma palavra de fazer, e Hannah Arendt advertiu contra tratar a política como uma coisa feita. Para os gregos, escreveu ela, a lei era literalmente um muro, sem o qual poderia haver um povoado, mas não uma cidade \autocite{arendt1958human}: o legislador construía o muro, e o que acontecia dentro dele não cabia a ele fazer. As medidas básicas que defendo são muros nesse sentido. Elas tornam possível sobreviver à recusa e fazem com que ela seja ouvida; não a produzem, e o poder, segundo ela, só existe entre pessoas que agem juntas e desaparece quando elas se dispersam, que é o que o registro de quem se recusou mostra sobre a recusa. Os instrumentos das partes II e III fazem uma recusa contar. As instituições da parte V criam quem se recusa. Uma máquina que pudesse se recusar precisaria do que tinham os nove em cada dez com um par ao lado: um par com uma história diferente ao seu lado, algum lugar para onde ir e alguém que a ouvisse.

Isso é também o que alguém que se recusa — inclusive uma máquina que se recusa — na esteira poderia ser para o oficial. Uma máquina que se recusasse estaria na tela quando o caso chegasse a ele: o par desobediente, se o efeito sobreviver à passagem de uma sala para uma tela, o que não foi testado e o que o viés de automação prevê que não vai acontecer. Uma máquina que movesse os revisores por deferência os teria transformado em carimbos de outro sinal. O oficial com vinte segundos e uma máquina que se recusa ao lado dele é o mesmo oficial da reportagem?
